% !TeX root = ../数值分析与科学计算笔记.tex
% 由主文档通过 \input 载入；请编译主文档。
\setcounter{chapter}{6} % 保留原书章号。
\chapter{矩阵特征值的数值计算}
\label{chap:na-eigen}

\section{特征对、归一化与残差}
求特征对要求 $Av=\lambda v$ 且 $v\ne0$。特征向量可以任意非零倍缩放，因此通常约定 $\|v\|_2=1$，比较时还要考虑符号或复相位。近似特征对的残差为 $r=Av-\mu v$。单个特征向量在重复或聚集特征值附近可能很敏感，此时不变子空间比特定基向量更有意义。

\begin{proposition}[对称矩阵的残差界]
设 $A$ 实对称，$\|v\|_2=1$，$\mu\in\mathbb R$。则至少有一个特征值满足 $|\lambda-\mu|\le\|Av-\mu v\|_2$。
\end{proposition}
\begin{proof}
取正交特征基 $q_i$，写 $v=\sum_i c_iq_i$，$\sum_i c_i^2=1$。则残差平方为 $\sum_i(\lambda_i-\mu)^2c_i^2$，不小于 $\min_i|\lambda_i-\mu|^2$，开平方即得结论。
\end{proof}
这一定理给出某个特征值接近的保证，并不直接保证 $v$ 接近某个唯一特征向量；后者还需要谱间隙。

\section{幂法与反幂法}
幂法输入非零 $v_0$，每步计算 $w_k=Av_k$，若 $w_k\ne0$ 则令 $v_{k+1}=w_k/\|w_k\|_2$。若 $A$ 可对角化，存在唯一模最大的特征值 $|\lambda_1|>|\lambda_2|\ge\cdots$，且 $v_0$ 在其特征向量方向的展开系数不为零，则归一化方向趋近于该主特征方向，误差的主导速率由 $|\lambda_2/\lambda_1|^k$ 控制。若主特征值为负，向量符号可以交替；严谨的收敛对象是方向。证明直接来自 $A^kv_0=\sum_i c_i\lambda_i^kv_i$ 并除以 $\lambda_1^k$。

对实对称矩阵，可用\emph{Rayleigh 商（Rayleigh quotient）}
\[
 \mu(v)=\frac{v^{\mathsf T}Av}{v^{\mathsf T}v}
\]
估计特征值。若单位向量与某特征向量夹角为 $\theta$，正交展开显示该商的误差为 $O(\sin^2\theta)$，常数受其他特征值距离控制，故其精度常优于方向误差的一阶量。

\begin{example}[缺失主方向的初值]
对 $A=\operatorname{diag}(3,1)$，初值 $(1,1)^{\mathsf T}$ 产生分量比 $3^k:1$，归一化方向趋于第一坐标轴；初值 $(0,1)^{\mathsf T}$ 却永远留在第二坐标轴。主特征值存在且简单并不足以保证任意非零初值都成功。
\end{example}

为寻找最接近给定移位 $\sigma$ 的特征值，反幂法每步解 $(A-\sigma I)w_k=v_k$ 再归一化。若 $\sigma$ 不是特征值，则逆矩阵的特征值为 $(\lambda_i-\sigma)^{-1}$；唯一最近且被初值激活的方向成为主方向。固定移位可复用矩阵分解；过于接近特征值时求解系统病态，应同时监测原特征问题的残差。

\section{QR 迭代与实际计算范围}
不带移位的 QR 迭代从 $A_0=A$ 出发，每步作 $A_k=Q_kR_k$，再令 $A_{k+1}=R_kQ_k=Q_k^{\mathsf T}A_kQ_k$。每一步是正交相似变换，保持特征值。带移位时分解 $A_k-\sigma_kI=Q_kR_k$，再取 $A_{k+1}=R_kQ_k+\sigma_kI$，同样保持相似性。

实际稠密算法先把一般矩阵化为 Hessenberg 形，对称矩阵化为三对角形，再使用移位与消去已收敛块的策略。一般实矩阵可能含复共轭特征值，因此实 Schur 形允许 $2\times2$ 对角块，不能要求所有实矩阵都迭代到实对角阵。本节说明算法结构，不给出一般带移位 QR 算法的完整收敛理论；课程实验宜先选实对称矩阵，并以残差和正交性核验结果。

\section{习题}
\begin{enumerate}[label=\arabic*.]
\item\label{exr:na-eigen-power} 对 $\operatorname{diag}(4,2,1)$ 及初值 $(1,1,1)^{\mathsf T}$ 写出第 $k$ 次归一化幂迭代向量，并指出主导方向误差比例。
\item\label{exr:na-eigen-rayleigh} 对 $A=\operatorname{diag}(\lambda_1,\lambda_2)$、$v=(\cos\theta,\sin\theta)^{\mathsf T}$，求 Rayleigh 商误差的精确表达式。
\item\label{exr:na-eigen-shift} 对 $A=\operatorname{diag}(1,3,8)$，移位 $\sigma=2.8$ 的反幂法通常趋向哪一个特征方向？注明初值条件。
\item\label{exr:na-eigen-rotation} 解释实旋转矩阵 $\begin{pmatrix}0&-1\\1&0\end{pmatrix}$ 为什么不能正交相似于实对角阵。
\end{enumerate}

